\section{Nested Classes}
Defining a class within another class/method, to group logically relevant
classes within the same encapsulation.\\
\textcolor{Bittersweet}{NOTE:\@ Container class CAN access private fields of
nested class!}
\begin{verbatim}
public/[default] class OuterClass {
   static class StaticNestedClass { ... }
   class InnerClass { ... }
}
\end{verbatim}
\paragraph{Access Modifiers}
Nested classes can be \verb$private$, \verb$protected$, \verb$[default]$,
\verb$public$. Private nested classes' instances can still be returned and
exist outside of the class, but their methods and fields cannot be used, even
if they are \verb$public$.\\
Should only be exposed $\iff$ no implementation details leaked, and is an essential part of
interface of outer class.
\subsection{Static Nested Class}
Associated with the containing \emph{class}. Can only access \textcolor{Blue}{static
fields and methods} of containing class.
\begin{verbatim}
Outer.Inner iObj = new Outer.Inner();
\end{verbatim}
\subsection{Inner Class}
Associated with the containing \emph{instance}. Can access \textcolor{Blue}{all
fields and methods} of containing class.
\begin{verbatim}
Outer oObj = new Outer();
Outer.Inner iObj = oObj.new Inner();
\end{verbatim}
\paragraph{Local Class}
Class defined within a function, scoped within the method. Can access
\textcolor{Blue}{fields and methods} of containing class, and
\textcolor{Blue}{local variables} of enclosing method. A local class in a
static method can only access static members of the containing class.
\paragraph{Anonymous Class}
Declaration AND instantiation of a local class in a single statement.
\begin{center}
  \verb$new SuperClass(args) { body }$
\end{center}
Fields, extra methods, instance initialisers and local classes are allowed in
\verb$body$, like a normal class \textcolor{Bittersweet}{without constructors}.
\subsection{Variable Capture}
Creating a copy of local variables inside the local class.\ \emph{Language
design decision} to only allow variables that are \verb$final$
\textcolor{Bittersweet}{or ``effectively final''}. Mutation not by
reassignment is still allowed.
\paragraph{Shadowing}
When the inner class declares a variable with the same name as one of the
enclosing class, it shadows the variable.\\
\verb$this.x$ refers to the inner class' member, while \verb$Outer.this.x$
refers to the outer class' member.
\paragraph{Stack and Heap}
On the stack and heap, a new instance of this class will include a copy of the
variables required in the inner class, \verb$Outer.this$, and the inner class'
variables.
\subsection{Lambda Expression}
\begin{center}
  \verb$@FunctionalInterface$
\end{center}
for interfaces with exactly \textbf{one} abstract method. Able to use lambda
expressions as syntactic sugar to simplify \emph{usage of functions as
first-class citizens}.\\
Similar to anonymous classes without introducing a new level of scoping.
Declarations, and variable references are \emph{interpreted just as they are}
in the enclosing environment. No shadowing issues \verb$:)$
\textcolor{Bittersweet}{Not allowed to redeclare variables, even as a parameter.}
\paragraph{Curried Functions}
\begin{center}
  $(X,Y)\rightarrow Z\equiv X\rightarrow Y\rightarrow Z$
\end{center}
Translating a $n$-ary function to a sequence of $n$ unary functions–a
higher-order function. Each lambda expression is a \textcolor{Blue}{closure},
\emph{storing data from the environment} where it is defined.
\paragraph{Method Referencing}
Referencing an existing method rather than specifying a new lambda. Can be used for:
\begin{itemize}
  \item Static method \verb$A::foo$ $\Rightarrow$ \verb$(args) -> A.foo(args)$
  \item Instance method\\
    \verb$aObj::foo$ $\Rightarrow$ \verb$args -> aObj.foo(args)$\\
    \verb$A::foo$$\Rightarrow$\verb$(args[1..])->args[0].foo(args[1..])$
  \item Constructors \verb$A::new$ $\Rightarrow$ \verb$args -> new A(args)$
\end{itemize}
During compilation, type inference is used to match the given reference to a
matching method.\ \textcolor{Bittersweet}{Compilation Error} if multiple matches
or ambiguity.
